% --- INICIO DE PLANTILLA PERSONALIZADA ---
\documentclass[oneside,11pt]{Latex/Classes/PhDthesisPSnPDF}

% --- Paquetes Originales ---
\usepackage{blindtext}
\usepackage{actuarialangle} 
\usepackage{tabularx}
\usepackage{float}
\usepackage{multirow, array}
\usepackage[usenames,dvipsnames,svgnames,table]{xcolor}
\usepackage{longtable}
\usepackage{latexsym,amsmath,amssymb,amsfonts}
\usepackage{enumitem}
\usepackage{xparse}
\usepackage{booktabs}
\usepackage[round, sort, numbers]{natbib}  
\usepackage{adjustbox}
\usepackage[utf8]{inputenc}
\usepackage{caption}
\usepackage{graphicx}
\usepackage{hyperref}
\usepackage{cleveref}

% Definición de colores
\definecolor{miblue}{RGB}{68, 114, 196}

% --- CAMBIO CRÍTICO 1: Usar \input para macros en el preámbulo ---
% \include da problemas antes del begin document. \input es seguro.
\input{Latex/Macros/MacroFile1}

% --- Definiciones de seguridad (por si la clase falla al cargar alguna variable) ---
\providecommand{\facultad}[1]{\def\lafacultad{#1}}
\providecommand{\escudofacultad}[1]{\def\elescudo{#1}}
\providecommand{\degree}[1]{\def\elgrado{#1}}
\providecommand{\director}[1]{\def\eldirector{#1}}
\providecommand{\lugar}[1]{\def\ellugar{#1}}
\providecommand{\degreedate}[1]{\def\lafecha{#1}}

% --- Configuración para bloques de código de R (Pandoc) ---

%%%%%%%%%%%%%%%%%%%%%%%%%%%%%%%%%%%%%%%%%%%%%%%%%%%%%%%%%%%%%%%%%%%%%%%%%%%%%%%%
%                         DATOS (Conectados a RMarkdown)                       %
%%%%%%%%%%%%%%%%%%%%%%%%%%%%%%%%%%%%%%%%%%%%%%%%%%%%%%%%%%%%%%%%%%%%%%%%%%%%%%%%
\title{Título del Trabajo Final}
\author{Eymi Borjas \\ Fabiola Bocaney}

% Lógica para Facultad: Si la defines en el YAML la usa, si no, usa la default

\facultad{Facultad de Ciencias Económicas y Sociales \\ Escuela de
Estadistica y Ciencias Actuariales}
  \escudofacultad{Latex/Classes/Escudos/faces}

\degree{Computación I}
\director{Jesus Ochoa \\ Oliver Triveño}
\degreedate{Marzo 2026} 
\lugar{Caracas}

\portadatrue 

% Metadatos PDF
\keywords{}
\subject{}

% --- CORRECCIÓN PARA PANDOC NUEVO (Imágenes) ---
\makeatletter
\providecommand{\pandocbounded}[1]{#1}
\makeatother


%%%%%%%%%%%%%%%%%%%%%%%%%%%%%%%%%%%%%%%%%%%%%%%%%%%%%
%                   DOCUMENTO                       %
%%%%%%%%%%%%%%%%%%%%%%%%%%%%%%%%%%%%%%%%%%%%%%%%%%%%%
\begin{document}

\maketitle

%%%%%%%%%%%%%%%%%%%%%%%%%%%%%%%%%%%%%%%%%%%%%%%%%%%%%
%                  PRÓLOGO                          %
%%%%%%%%%%%%%%%%%%%%%%%%%%%%%%%%%%%%%%%%%%%%%%%%%%%%%
\frontmatter

% Puedes agregar dedicatorias desde el YAML si quieres, o dejarlas fijas aquí

%%%%%%%%%%%%%%%%%%%%%%%%%%%%%%%%%%%%%%%%%%%%%%%%%%%%%
%                   ÍNDICES                         %
%%%%%%%%%%%%%%%%%%%%%%%%%%%%%%%%%%%%%%%%%%%%%%%%%%%%%
\setcounter{secnumdepth}{3}
\setcounter{tocdepth}{3}

\tableofcontents



\cleardoublepage

  \cleardoublepage       % Empieza en página derecha (estándar en tesis)
  \chapter*{Resumen}     % Crea el título "Resumen" sin numeración (Capítulo *)
  \addcontentsline{toc}{chapter}{Resumen} % (Opcional) Lo añade al índice
  
  Resumen del trabajo de
investigación\ldots{}             % Aquí se pega el texto que escribiste en el YAML
  
  \cleardoublepage

%%%%%%%%%%%%%%%%%%%%%%%%%%%%%%%%%%%%%%%%%%%%%%%%%%%%%
%                   CONTENIDO (El Corazón)          %
%%%%%%%%%%%%%%%%%%%%%%%%%%%%%%%%%%%%%%%%%%%%%%%%%%%%%
\mainmatter
\def\baselinestretch{1.5}

% --- CAMBIO CRÍTICO 2: Aquí RMarkdown inyecta todo tu texto ---
\chapter*{Introducción}

Este documento técnico se encuentra estructurado en cinco capítulos
fundamentales:

Capítulo I Se expone la relevancia de analizar la consistencia deportiva
frente a la efectividad aislada, formulando las preguntas de
investigación y delimitando el alcance.

Capítulo II Establece el Marco Teórico, definiendo los conceptos
estadísticos clave como, necesarios para la interpretación de los datos.

Capítulo III Describe el Marco Metodológico, detallando el proceso de
limpieza de datos (ETL), las pruebas de hipótesis y los algoritmos
utilizados para validar la calidad de la información.

Capítulo IV Presenta el Análisis de Resultados, donde se exhiben las
tablas y gráficos generados dinámicamente.

Capítulo V Expone los hallazgos finales sobre qué factores determinaron
la victoria en 2022. Se concluye sobre la estabilidad de los equipos y
se ofrecen recomendaciones estratégicas basadas en la evidencia
estadística obtenida.

\chapter{Planteamiento del Problema}

En el baloncesto de élite, ganar un partido no siempre depende de qué
tan alto sea el promedio de un equipo, sino de qué tan consistente pueda
ser ese rendimiento bajo presión. El análisis de los 542 registros de la
NBA en 2022 revela una interrogante crítica: ¿se pierden los partidos
por una falta de talento general o por una irregularidad drástica en
momentos clave?

El problema radica en que el porcentaje de tiro (FG\_PCT) suele ser
volátil. Cuando un equipo pierde, no siempre queda claro si es por una
ineficiencia constante o por una alta variabilidad (inestabilidad) en
sus lanzamientos. Además, existe una brecha en la comprensión de cómo
factores específicos, como superar el 50\% de efectividad de campo o
dominar la batalla de rebotes, actúan como predictores determinantes de
la victoria.\\

Argumento final. Debido a esto, se plantea las siguientes preguntas de
investigación:\\

\textbf{1.} ¿Es la derrota en la NBA 2022 consecuencia de una
ineficiencia constante en el tiro o se debe a una alta volatilidad e
inestabilidad en el porcentaje de campo (FG\_PCT)?\\

\textbf{2.} ¿En qué medida el superar el umbral del 50\% de efectividad
en tiros de campo o ganar la batalla de rebotes garantiza
estadísticamente una victoria en el marcador final?\\

\textbf{3.}¿Cómo se transforma el volumen de puntos anotados (PTS) a
medida que un equipo escala en sus rangos de eficiencia de tiro?\\

\section{Justificación}

El presente trabajo se delimita al análisis de la temporada 2022 con el
fin de garantizar la homogeneidad en el contexto táctico y reglamentario
de los datos. La selección de este periodo y de las variables
correspondientes responde a los siguientes criterios:

Se priorizan aquellos indicadores que reflejan directamente la
eficiencia ofensiva y el control del juego, permitiendo un análisis con
mayor valor interpretativo. El uso de este subconjunto de datos facilita
la aplicación de herramientas de tendencia central, dispersión y tablas
bivariantes, objetivos principales de este estudio.

Al centrarse en un año específico, se eliminan ruidos estadísticos
provocados por la evolución histórica del juego, permitiendo
conclusiones más precisas sobre el rendimiento contemporáneo en la
NBA.\\

\section{Objetivos}

\subsection{Objetivo General}
\begin{itemize}
  \item{Evaluar la incidencia de la estabilidad en el tiro y la relación de variables de efectividad y posesión en el resultado final de los partidos de la NBA en el año 2022, mediante el análisis de medidas de dispersión y tablas de contingencia.}
\end{itemize}

\subsection{Objetivos Específicos}
\begin{itemize}
  \item{Analizar el porcentaje de tiro $(FG_PCT)$ en las derrotas para determinar si los partidos se pierden por una baja eficiencia constante o por una alta inestabilidad en el tiro.}
  \item{Categorizar el porcentaje de tiro $(FG_PCT)$ en niveles para observar cómo cambia el promedio de puntos anotados en cada nivel.}
  \item{Construir tablas para cuantificar la frecuencia de victorias}
 \item{Generar un reporte técnico reproducible en formato PDF que integre visualizaciones estadísticas y tablas resumen utilizando RMarkdown y LaTeX.}
\end{itemize}

\section{Cobertura}

\subsection{Horizontal}

Describir la cobertura Horizontal del trabajo de investigación\\

\begin{table}[ht]
\centering
\caption{Variables consideradas} 
\resizebox{10cm}{!} {
\begin{tabular}{ccc}
  \hline
\ \ \ \ \ NOMBRE DE LA VARIABLE \\ 
  \hline
   VARIABLE 1 \\ 
VARIABLE 2\\ 
VARIABLE n\\ 
   \hline
\end{tabular}
}
\end{table}

\subsection{Vertical}

Describir la cobertura Vertical del trabajo de investigación

\section{Periodo de Referencia}

Describir la fuente y el Periodo de Referencia utilizado en el trabajo
de investigación

\section{Variables}

Describir la variable en estudio en el trabajo de investigación

\chapter{Marco Teórico}

\section{Bases Teóricas}

Resumen de las bases teoricas explicadas en el capitulo y su
importancia\\

\subsection{Definición 1}

Descripción de la definición 1\\

\subsection{Definición 2}

Descripción de la definición 2\\

\subsection{Definición n}

Descripción de la definición n\\

\chapter{Marco Metódico}

En esta sección se presentan las metodologías y/o test necesarios
relacionados con la práctica divulgada en el marco teórico, así como las
fuentes de datos.

\section{Bases metódicas utilizadas en el trabajo de investigación}

Desarrollo\\

\subsection{Procesamiento y limpieza de datos (ETL}

\begin{verbatim}
## Rows: 542
## Columns: 14
## $ SEASON         <chr> "2022", "2022", "2022", "2022", "2022", "2022", "2022",~
## $ PTS_home       <dbl> 126, 120, 114, 113, 108, 112, 143, 106, 110, 99, 101, 1~
## $ FG_PCT_home    <dbl> 0.484, 0.488, 0.482, 0.441, 0.429, 0.386, 0.643, 0.553,~
## $ FT_PCT_home    <dbl> 0.926, 0.952, 0.786, 0.909, 1.000, 0.840, 0.875, 0.611,~
## $ FG3_PCT_home   <dbl> 0.382, 0.457, 0.313, 0.297, 0.378, 0.317, 0.636, 0.423,~
## $ AST_home       <dbl> 25, 16, 22, 27, 22, 26, 42, 25, 22, 23, 19, 29, 29, 22,~
## $ REB_home       <dbl> 46, 40, 37, 49, 47, 62, 32, 38, 49, 39, 37, 46, 48, 43,~
## $ PTS_away       <dbl> 117, 112, 106, 93, 110, 117, 113, 113, 116, 104, 98, 12~
## $ FG_PCT_away    <dbl> 0.478, 0.561, 0.470, 0.392, 0.500, 0.469, 0.494, 0.447,~
## $ FT_PCT_away    <dbl> 0.815, 0.765, 0.682, 0.735, 0.773, 0.778, 0.760, 0.909,~
## $ FG3_PCT_away   <dbl> 0.321, 0.333, 0.433, 0.261, 0.292, 0.462, 0.364, 0.265,~
## $ AST_away       <dbl> 23, 20, 20, 15, 20, 27, 32, 17, 19, 17, 29, 25, 25, 27,~
## $ REB_away       <dbl> 44, 37, 46, 46, 47, 47, 36, 38, 45, 39, 36, 39, 40, 43,~
## $ HOME_TEAM_WINS <fct> 1, 1, 1, 1, 0, 0, 1, 0, 0, 0, 1, 1, 1, 0, 0, 1, 0, 1, 1~
\end{verbatim}

Descripción\\

\subsection{Item 2}

Descripción\\

\subsubsection{Item 2.2}

Descripción\\

\subsection{Item n}

Descripción\\

\chapter{Análisis de Resultados}

\{\small Este capítulo presenta los resultados obtenidos al aplicar las
metodologías descritas en el capítulo anterior, así como del análisis
estadístico descriptivo de las variables estudiadas y las ventajas y
desventajas al aplicar dichos tests.\}

\section{Presentación de resultados 1}

Descripción\\

\section{Presentación de resultados 2}

Descripción\\

\section{Presentación de resultados n}

Descripción\\

\chapter{Conclusiones y Recomendaciones}

\textbf{1.-} Conclusion 1 del trabajo de investigación.\\

\textbf{Recomendación} Recomendación 1 del trabajo de investigación.\\

\textbf{2.-} Conclusion 2 del trabajo de investigación.\\

\textbf{Recomendación} Recomendación 2 del trabajo de investigación.\\

\textbf{n.-} Conclusion n del trabajo de investigación.\\

\textbf{Recomendación} Recomendación n del trabajo de investigación.\\

\appendix
\renewcommand{\thechapter}{\Alph{chapter}}

\chapter{Anexo A}

\chapter{Anexo B}

\chapter{Anexo C}

%%%%%%%%%%%%%%%%%%%%%%%%%%%%%%%%%%%%%%%%%%%%%%%%%%%%%
%                   REFERENCIAS                     %
%%%%%%%%%%%%%%%%%%%%%%%%%%%%%%%%%%%%%%%%%%%%%%%%%%%%%
% Mantenemos la estructura natbib de tu plantilla
\renewcommand{\bibname}{Lista de Referencias}
\bibliographystyle{apalike} 
\bibliography{bibliografia.bib} 

%%%%%%%%%%%%%%%%%%%%%%%%%%%%%%%%%%%%%%%%%%%%%%%%%%%%%
%                   APÉNDICES                       %
%%%%%%%%%%%%%%%%%%%%%%%%%%%%%%%%%%%%%%%%%%%%%%%%%%%%%

\end{document}
